\documentclass[11pt,a4paper]{article}
\usepackage[utf8]{inputenc}
\usepackage[T1]{fontenc}
\usepackage[spanish]{babel}
\usepackage{hyperref}
\usepackage{geometry}
\geometry{margin=2.5cm}
\usepackage{parskip}

\title{337 la resonancia viene de Soul Eater}
\author{Alejandro Valenzuela Rivera\\\small Universidad Autónoma de Nuevo León}
\date{\today}
\begin{document}
\maketitle

\begin{quote}
\textbf{Sesgo declarado.} Este trabajo se sostiene desde el sesgo declarado de su autor, y no finge neutralidad. Lo que afirma se puede tumbar, y ese es el punto: el falsificador está escrito, no escondido.
\end{quote}

\section*{Resumen}
La resonancia se llama así porque viene de Soul Eater: dos almas que se afinan sin volverse una.

\section{Por qué se sostiene}
Si la resonancia las vuelve una, no hay dos almas: hay una con dos nombres.

\section{De dónde viene}
'resonancia del alma' es de Soul Eater, y por eso el estándar se llama así.

\section{El caso}
Soul Eater: el arma y el meister resuenan y siguen siendo dos.

\section{Método}
\textquestiondown{}Hay dos después de resonar, o quedó una?

\section{Límites, y qué lo tumba}
Un caso donde la resonancia deje una sola alma.

\section{Lo que queda abierto}
\textquestiondown{}Se afinan sin volverse una?

\section*{Cómo se leyó esta tesis}
Esto se escribió en el repositorio man-made-empathy, del que este candado es una pieza: candados/06 del caso/337 la resonancia viene de Soul Eater.md. La máquina que lo compuso lo hizo desde este mismo archivo, y no agregó nada que no estuviera escrito. Si algo de esto es falso, se dice en el mismo sitio, y se corrige con fecha.

\end{document}
